\keys{educational data mining; student performance prediction; machine learning; pedagogical risk; learning analytics; tabular temporal modeling.} % Palavras chave em inglês separadas por ponto e vírgula


% Não deixe linha em branco após \begin{abstract}{Title in English}
\begin{abstract}{Predictive system for school performance to support pedagogical decision-making}
This undergraduate thesis presents the development of a predictive system for school environments. Its goal is to anticipate academic difficulties and support decision-making by teachers, pedagogical coordinators, and administrative staff. The research problem investigates whether historical school data, such as grades, absences, payments, and academic records, can be integrated into an analytical architecture capable of predicting a student's next grade and identifying pedagogical risk signals before failure or severe performance decline occurs. The project initially explored a graph-based approach, but evolved into a tabular temporal strategy after the analyses indicated stronger adherence between the problem and supervised models applied to longitudinal data. The final solution was organized around the \texttt{school\_predictor} package, with a private local layer for database preparation and extraction, a public CSV-based interface, feature engineering, temporal validation, model selection, error analysis, and report generation. Baselines, Random Forest, and HistGradientBoosting were compared for both next-grade regression and classification of risk for grades below 6.0. For numeric next-grade prediction, the best setting used a minimum history of one prior observation per subject and achieved MAE of 0.7659. For risk alerting, the best setting used a minimum history of two prior observations, achieving F1 of 0.7436 and recall of 0.8794 in temporal testing. In addition to predictive modeling, a school reporting layer was developed with outputs for teachers, coordination, and school administration, including an executive ranking of priority students. The study concludes that a temporal analytical architecture oriented toward institutional use, combined with interpretation, prioritization, and data protection mechanisms, can support earlier pedagogical intervention.
\end{abstract}
